\documentclass[10pt]{article}
\usepackage[T1]{fontenc}
\usepackage[utf8]{inputenc}
\usepackage[french]{babel}
\usepackage{amsmath,amssymb,mathrsfs,mathpazo}
\usepackage{geometry}
\geometry{a4paper,margin=2cm}
\usepackage{graphicx,array,tabularx,tabularray,colortbl,multirow,multicol,booktabs,enumitem,pifont,ulem,fancybox,diagbox,cancel}
\usepackage{pstricks,pstricks-add,pst-plot,pst-tree,pst-node,pst-text,pst-func,pst-math,pst-bspline,pst-3dplot,pst-eucl,pst-solides3d}
\let\pscalculate\undefined
\usepackage{pst-bezier}
\usepackage{tikz,pgfplots,tkz-tab}
\usetikzlibrary{arrows,arrows.meta,patterns,calc,positioning,trees,decorations.markings,shapes.geometric,shapes.misc,shapes.arrows,intersections,decorations.pathmorphing,decorations.pathreplacing,backgrounds,shadows,babel}
\pgfplotsset{compat=1.18}
\usepgfplotslibrary{fillbetween}
\usepackage[np]{numprint}
\usepackage{listings,makecell,pas-tableur,twoopt}
\let\barre\undefined
\usepackage[pstricks]{tablvar}
\definecolor{darkgreen}{RGB}{0,100,0}
\providecommand{\red}{\color{red}}
\providecommand{\blue}{\color{blue}}
\providecommand{\green}{\color{green}}
\providecommand{\black}{\color{black}}
\providecommand{\white}{\color{white}}
\providecommand{\gray}{\color{gray}}
\providecommand{\cyan}{\color{cyan}}
\providecommand{\magenta}{\color{magenta}}
\providecommand{\yellow}{\color{yellow}}
\providecommand{\violet}{\color{violet}}
\providecommand{\orange}{\color{orange}}
\providecommand{\euro}{\texteuro}
\providecommand{\EUR}{\texteuro}
\providecommand{\ieme}{\textsuperscript{e}}
\providecommand{\degres}{\ensuremath{^\circ}}
\providecommand{\degree}{\ensuremath{^\circ}}
\providecommand{\vv}[1]{\overrightarrow{#1}}
\providecommand{\Coord}[1]{#1}
\providecommand{\decosix}{}
\providecommand{\starredbullet}{\ensuremath{\bullet}}
\pagestyle{empty}
\setlength{\parindent}{0pt}
\usepackage{adjustbox}\DeclareUnicodeCharacter{2236}{:}\DeclareUnicodeCharacter{2190}{\ensuremath{\leftarrow}}
\providecommand{\R}{\mathbb{R}}
\providecommand{\N}{\mathbb{N}}
\providecommand{\Z}{\mathbb{Z}}
\providecommand{\Q}{\mathbb{Q}}
\providecommand{\D}{\mathbb{D}}
\providecommand{\C}{\mathbb{C}}
\providecommand{\e}{{\mathrm{e}}}
\providecommand{\vect}[1]{\overrightarrow{#1}}
\providecommand{\vectt}[1]{\overrightarrow{\mathrm{#1}}}
\providecommand{\vvt}[1]{\overrightarrow{\mathrm{#1}}}
\providecommand{\np}[2][]{{#2~#1}}
\providecommand{\vv}[1]{{\overrightarrow{#1}}}
\renewcommand{\d}{{\mathrm{d}}}
\providecommand{\degres}{{^\circ}}
\providecommand{\degre}{{^\circ}}
\providecommand{\blue}{\color{blue}}
\providecommand{\red}{\color{red}}
\providecommand{\enskip}{\quad}
\providecommand{\llbracket}{\lbrack\!\lbrack}
\providecommand{\rrbracket}{\rbrack\!\rbrack}
\providecommand{\euro}{{\text{€}}}
\providecommand{\no}{{\text{n°}}}
\providecommand{\interval}[2]{{[#1;#2]}}
\providecommand{\intervFO}[2]{{[#1;#2[}}
\providecommand{\dotfill}{\ldots}

\providecommand{\ds}{\displaystyle}
\providecommand{\textding}[1]{#1}
\providecommand{\Oij}{\(\mathrm{O};\vec{\imath},\vec{\jmath}\)}
\providecommand{\Oijk}{\(\mathrm{O};\vec{\imath},\vec{\jmath},\vec{k}\)}
\providecommand{\Ouv}{\(\mathrm{O};\vec{u},\vec{v}\)}
\providecommand{\intervFF}[2]{\ensuremath{[#1;#2]}}
\renewcommand{\euro}{\ifmmode\text{\texteuro}\else\texteuro\fi}
\DeclareUnicodeCharacter{00B0}{\ifmmode^{\circ}\else\textdegree\fi}
\DeclareUnicodeCharacter{20AC}{\ifmmode\text{\texteuro}\else\texteuro\fi}
% Énoncés : Éducation nationale. Transcription et corrigés : APMEP, Denis Vergès et contributeurs. Découpage et classement : CoopMaths, Sébastien Lozano, Éric Elter et contributeurs. Adaptation de la présentation : Maths Explorer.
% CC BY-SA : https://coopmaths.fr/www/about/
% https://coopmaths.fr/alea/static/bac/2022/tex/bac_2022_09_sujet1_polynesie_4.tex
\begin{document}
\noindent\textit{APMEP / CoopMaths. CC BY-SA. Présentation adaptée par Maths Explorer.}\par\medskip
\medskip

On considère le cube ABCDEFGH.

On note I le milieu du segment[EH] et on considère le triangle CFI.

L'espace est muni du repère orthonormé $\left(\text{A}~;~ \vect{\text{AB}},\: \vect{\text{AD}},\: \vect{\text{AE}}\right)$ et on admet que le point I a
pour coordonnées $\left(0~;~\dfrac12;~1\right)$ dans ce repère.

\begin{center}
\psset{unit=1cm}
\begin{pspicture}(7.4,8)
\pspolygon(0.2,0.8)(4.6,0.2)(4.6,4.6)(0.2,5.2)
\psline(4.6,0.2)(7.2,2.6)(7.2,7)(4.6,4.6)
\psline(7.2,7)(2.8,7.6)(0.2,5.2)
\psline[linestyle=dashed](0.2,0.8)(2.8,3.2)(7.2,2.6)
\psline[linestyle=dashed](2.8,3.2)(2.8,7.6)
\psline(5,7.3)(0.2,5.2)(4.6,0.2)
\psline[linestyle=dashed](5,7.3)(4.6,0.2)
\psdots(2.8,3.2)(0.2,0.8)(4.6,0.2)(7.2,2.6)(2.8,7.6)(0.2,5.2)(4.6,4.6)(7.2,7)(5,7.3)(2.4,4.9)
\uput[l](2.8,3.2){A} \uput[l](0.2,0.8){B} \uput[dr](4.6,0.2){C} \uput[r](7.2,2.6){D}
\uput[u](2.8,7.6){E} \uput[l](0.2,5.2){F} \uput[dl](4.6,4.6){G} \uput[ur](7.2,7){H}
\uput[u](5,7.3){I} \uput[u](2.4,4.9){J}
\end{pspicture}
\end{center}
\smallskip

\begin{enumerate}
\item 
	\begin{enumerate}
		\item Donner sans justifier les coordonnées des points C, F et G.
		\item Démontrer que le vecteur $\vect{n}\begin{pmatrix}1\\2\\2\end{pmatrix}$ est normal au plan (CFI).
		\item Vérifier qu'une équation cartésienne du plan (CFI) est : $x + 2y + 2z - 3 = 0$.
	\end{enumerate}	
\item  On note $d$ la droite passant par G et orthogonale au plan (CFI).
	\begin{enumerate}
		\item Déterminer une représentation paramétrique de la droite $d$.
		\item Démontrer que le point K$\left(\dfrac79~;~\dfrac59~;~\dfrac59\right)$ est le projeté orthogonal du point G sur le plan (CFI).
		\item Déduire des questions précédentes que la distance du point G au plan (CFI) est
égale à $\dfrac23$.
	\end{enumerate}
\item On considère la pyramide GCFI.

\emph{On rappelle que le volume $V$ d'une pyramide est donné par la formule} 

\[V = \frac13 \times b \times h,\]

\emph{où $b$ est l'aire d'une base et $h$ la hauteur associée à cette base}.
	\begin{enumerate}
		\item Démontrer que le volume de la pyramide GCFI est égal à $\dfrac16$, exprimé en unité de volume.
		\item En déduire l'aire du triangle CFI, en unité d'aire.
	\end{enumerate}
\end{enumerate}






\newcounter{INDICE} 
\newcommandtwoopt{\SuiteRec}[5][u][0]{
\getcoor 
\infixtoRPN{#4} 
\psplot{\xmin}{\xmax}{\RPN} 
\SpecialCoor 


\psline[linestyle=dashed](\xmin,\xmin)(\xmax,\xmax)  
\pstVerb{ /x #3 def } 
\psline[ArrowInside=->,ArrowInsideNo=1](! x 0 )(! x \RPN ) 
\psline[ArrowInside=->,ArrowInsideNo=1](! x \RPN )(! \RPN\space dup) 
\psline{-}(! x 0.03 )(! x 0.03 neg) 
\rput(!x 0.08 neg ){$#1_{#2}$} 
\pstVerb{ /x \RPN\space def} 
\setcounter{INDICE}{#2} 
\addtocounter{INDICE}{1} 
\multido{\I=\theINDICE+1}{#5}{
\psline{-}(! x 0.03 )(! x 0.03 neg) 
\rput(!x 0.08 neg ){$#1_{\I}$} 
\psline[linestyle=dotted](! x 0 )(! x x) 
\psline[ArrowInside=->,ArrowInsideNo=1](! x x )(! x \RPN ) 
\psline[ArrowInside=->,ArrowInsideNo=1](! x \RPN )(! \RPN\space dup) 
\pstVerb{ /x \RPN\space def} 
}
}
\newcommandtwoopt{\SuiteRecbis}[5][u][0]{
\getcoor 
\infixtoRPN{1.9 x mul 1 x sub mul} 
\SpecialCoor 
\psline[linestyle=dashed](\xmin,\xmin)(\xmax,\xmax)  
\pstVerb{ /x #3 def } 
\psline[ArrowInside=->,ArrowInsideNo=1](! x 0 )(! x \RPN ) 
\psline[ArrowInside=->,ArrowInsideNo=1](! x \RPN )(! \RPN\space dup) 
\psline{-}(! x 0.03 )(! x 0.03 neg) 
\rput(!x 0.08 neg ){$#1_{#2}$} 
\pstVerb{ /x \RPN\space def} 
\setcounter{INDICE}{#2} 
\addtocounter{INDICE}{1} 
\multido{\I=\theINDICE+1}{4}{
\psline{-}(! x 0.03 )(! x 0.03 neg) 
\rput(!x 0.08 neg ){$#1_{\I}$} 
\psline[linestyle=dotted](! x 0 )(! x x) 
\psline[ArrowInside=->,ArrowInsideNo=1](! x x )(! x \RPN ) 
\psline[ArrowInside=->,ArrowInsideNo=1](! x \RPN )(! \RPN\space dup) 
\pstVerb{ /x \RPN\space def} 
}
}

\end{document}
